\labeledsection{Methods}
Before starting with the experiment we first need to find out how to properly implement \textit{Dependency-Injection}.
For simplicity, we implement the \textit{plugin-system} within our application and a plugin will only perform a small task.
We also only do this in \textit{Java}, the same can be applied to almost any language but will be significantly different from the implementation.

\textit{Dependency-Injection} in \textit{Java} can be done in multiple ways.
One can do everything manually, define their own class loaders and load a \textit{plugin} into memory.
However, there is an official way of modifying the class-path:
The Java Agent API / \texttt{java.lang.Instrument} package.

From the \textcite{Instrument-JavaDoc}:
\begin{quote}
    [The package p]rovides services that allow Java programming language agents to instrument programs running on the JVM.\newline
    The mechanism for instrumentation is modification of the byte-codes of methods.
\end{quote}~\parencite{Instrument-JavaDoc}\\



\begin{enumerate}
    \item How does Java Dependency-Injection work?
            Implemented Plugin-System
            External dependencies
    \item experiment
\end{enumerate}